\section{Physics-Informed 3D Surface Reconstruction and Dynamic Suspension Simulation}
\label{sec:3d_reconstruction_and_simulation}

\subsection{Architectural Blueprint: From Monocular Frustum to Dynamic Response}
In the absence of direct synchronous hardware pairing between vehicle dashcams and internal chassis accelerometers, classical supervised end-to-end regression fails. To solve this uncoupled multi-modal inference challenge, we construct a physics-informed operational pipeline that couples monocular projective geometry with classical multi-body vehicle dynamics. 

The core architectural paradigm executes in four sequential stages:
\begin{enumerate}
    \item \textbf{Metric Unprojection and Frustum Mapping}: 2D bounding boxes detected across the image plane $(u, v)$ are mapped onto the Euclidean 3D road surface coordinate manifold $(X, Y, Z)$ using intrinsic camera pinhole geometry calibrated against ground plane mounting constraints.
    \item \textbf{3D RANSAC Road Reference Plane Fitting}: Pavement pixels situated in an annular peripheral ring around each candidate cavity are extracted to fit an unperturbed ground plane $Ax + By + Cz + D = 0$. Deviations beneath this reference plane yield metric depression depth profiles $z_{\text{pothole}}(x, y)$.
    \item \textbf{2-Degree-of-Freedom (2-DOF) Quarter-Car Dynamic ODE Simulation}: The reconstructed longitudinal cavity elevation profile $z_{\text{road}}(t) = z_{\text{pothole}}(v \cdot t)$ is injected into a quarter-car dynamic system as time-domain ground displacement, simulating the instantaneous vertical accelerations $\hat{a}_z(t)$ of the sprung and unsprung masses.
    \item \textbf{Empirical Dynamic Consistency Validation}: Simulated acceleration peaks and power spectra are benchmarked against our curated empirical library of real-world vehicular shock signatures, computing a physical confidence score $C_{\text{dyn}} \in [0.0, 1.0]$ and categorizing road distress into three standardized civil severity tiers.
\end{enumerate}

\begin{designreflection}{Why RANSAC Plane Fitting Outperforms Direct Depth Neural Networks}
When exploring 3D depth extraction from single monocular images, one could deploy zero-shot monocular depth estimation networks like Depth Anything or ZoeDepth. However, monocular depth networks suffer from an inherent projective ambiguity: scale-and-shift drift. While a neural network predicts relative depth with high perceptual fidelity, it cannot distinguish between a $3\,\text{cm}$ surface depression on an asphalt road and a $12\,\text{cm}$ deep crater without metric reference calibration.

To establish absolute metric ground truth, I designed an annular peripheral RANSAC plane fitting algorithm. By sampling the unperturbed pavement immediately enclosing the pothole cavity boundary, we construct a local planar datum. Because this local planar ring represents the true undisturbed road surface, any negative orthogonal offset represents actual physical cavity depth. By measuring relative displacement against this RANSAC datum rather than relying on ungrounded global depth predictions, our system eliminates neural scale drift and yields robust, sub-centimeter depression profiles.
\end{designreflection}

\subsection{Formulation of 3D Ground Plane Estimation and Metric Depth Profiling}

\subsubsection{Perspective Frustum Unprojection}
Consider an optical dashcam mounted at vertical elevation $H_{\text{cam}} = 1.25\,\text{m}$ above the flat road plane, oriented with optical axis parallel to the longitudinal vehicle motion vector. Let the image sensor have dimensions $W_i \times H_i$, horizontal field of view $\text{FOV}_h = 70^\circ$, and calibrated horizon ordinate $y_{\text{horizon}} = 0.45$. The normalized focal length is given by:

\begin{equation}
    f_{\text{norm}} = \frac{1}{2 \tan(\text{FOV}_h / 2)} \approx 0.714
    \label{eq:fnorm_recon}
\end{equation}

For any pixel coordinate $(u, v)$ in the lower image half ($v / H_i > y_{\text{horizon}}$), the metric longitudinal distance $Z$ ahead of the front bumper is dictated by perspective ground intersection:
\begin{equation}
    Z(u, v) = \frac{H_{\text{cam}}}{\frac{v}{H_i} - y_{\text{horizon}}} \quad [\text{meters}]
    \label{eq:z_distance_recon}
\end{equation}

The corresponding transverse lateral coordinate $X$ and vertical coordinate $Y$ in the camera frame are computed via:
\begin{equation}
    X(u, v) = \frac{\left(\frac{u}{W_i} - 0.5\right) \cdot Z(u, v)}{f_{\text{norm}}}, \quad Y(u, v) = \frac{\left(\frac{v}{H_i} - y_{\text{horizon}}\right) \cdot Z(u, v)}{f_{\text{norm}}}
    \label{eq:xy_coordinates_recon}
\end{equation}

\subsubsection{RANSAC Estimation of the Undisturbed Pavement Datum}
Given a detected pothole bounding box with pixel boundaries $[x_1, y_1, x_2, y_2]$, the interior pixels represent the disrupted road cavity. To estimate the undisturbed road plane, we define an annular peripheral boundary ring $\Omega_{\text{datum}}$:
\begin{equation}
    \Omega_{\text{datum}} = \left\{ (u, v) \mid u \in [x_1 - \delta_x, x_2 + \delta_x], \, v \in [y_1 - \delta_y, y_2 + \delta_y] \right\} \setminus [x_1, x_2] \times [y_1, y_2]
\end{equation}
where padding margins $\delta_x = 0.25(x_2 - x_1)$ and $\delta_y = 0.25(y_2 - y_1)$ ensure adequate sampling of unperturbed asphalt.

We sample $N_{\text{pts}} \ge 50$ 3D spatial points $\mathbf{P}_i = [X_i, Y_i, Z_i]^T \in \Omega_{\text{datum}}$ and fit a planar manifold $\mathbf{n} \cdot \mathbf{P} + D = 0$, where unit normal $\mathbf{n} = [A, B, C]^T$ ($\|\mathbf{n}\| = 1$) via Random Sample Consensus (RANSAC):
\begin{enumerate}
    \item Randomly select three non-collinear datum points $\mathbf{P}_1, \mathbf{P}_2, \mathbf{P}_3$.
    \item Compute candidate plane normal: $\mathbf{n}_{\text{cand}} = \frac{(\mathbf{P}_2 - \mathbf{P}_1) \times (\mathbf{P}_3 - \mathbf{P}_1)}{\|(\mathbf{P}_2 - \mathbf{P}_1) \times (\mathbf{P}_3 - \mathbf{P}_1)\|}$ and plane offset $D_{\text{cand}} = -\mathbf{n}_{\text{cand}} \cdot \mathbf{P}_1$.
    \item Evaluate orthogonal distance residuals for all candidate points:
    \begin{equation}
        d_i = \left| \mathbf{n}_{\text{cand}} \cdot \mathbf{P}_i + D_{\text{cand}} \right|
    \end{equation}
    \item A point is classified as an inlier if $d_i < \tau_{\text{plane}} = 0.03\,\text{m}$.
    \item Iterate for 200 trials and select the normal vector maximizing total inlier cardinality.
\end{enumerate}

\subsubsection{Extraction of Metric Depth, Aperture, and Volumetric Displacement}
With the reference plane $\mathbf{n} = [A, B, C]^T$ established, the ideal undisturbed road elevation $Z_{\text{plane}}(u, v)$ for any pixel inside the cavity is:
\begin{equation}
    Z_{\text{plane}}(u, v) = -\frac{A \cdot X(u, v) + B \cdot Y(u, v) + D}{C}
\end{equation}

The metric depression depth profile $d(u, v)$ (in centimeters, defined positive downwards into the road surface) is:
\begin{equation}
    d(u, v) = \max\left(0.0, \, \left( Z(u, v) - Z_{\text{plane}}(u, v) \right) \times 100 \right) \quad [\text{cm}]
    \label{eq:depression_profile}
\end{equation}

From this discrete 2D depression manifold, we extract three critical physical descriptors:
\begin{enumerate}
    \item \textbf{Maximum Metric Depth ($d_{\max}$)}:
    To reject camera sensor salt-and-pepper noise and specular reflections, $d_{\max}$ is evaluated as the 95th percentile of the depression distribution:
    \begin{equation}
        d_{\max} = \text{Percentile}_{95}\left( \{ d(u, v) \mid (u, v) \in \text{Cavity} \} \right) \quad [\text{cm}]
    \end{equation}
    
    \item \textbf{Physical Planar Dimensions ($W_{\text{phys}}, L_{\text{phys}}$)}:
    \begin{equation}
        W_{\text{phys}} = \frac{(x_2 - x_1) \cdot Z_{\text{center}}}{f_{\text{norm}} \cdot W_i} \times 100, \quad L_{\text{phys}} = \frac{(y_2 - y_1) \cdot Z_{\text{center}}}{f_{\text{norm}} \cdot H_i} \times 100 \quad [\text{cm}]
    \end{equation}
    
    \item \textbf{Volumetric Cavity Displacement ($V$)}:
    Computed via 2D numerical Riemann surface integration:
    \begin{equation}
        V = \sum_{u=x_1}^{x_2} \sum_{v=y_1}^{y_2} d(u, v) \cdot \Delta A_{\text{pixel}} \quad [\text{cm}^3]
    \end{equation}
    where $\Delta A_{\text{pixel}} = \left(\frac{W_{\text{phys}}}{x_2 - x_1}\right) \times \left(\frac{L_{\text{phys}}}{y_2 - y_1}\right)$.
\end{enumerate}

\subsection{Two-Degree-of-Freedom Quarter-Car Dynamic Model}

\subsubsection{Dynamic Governing Differential Equations}
To validate whether the reconstructed 3D pothole geometry produces physically plausible vertical shocks on passing vehicles, we implement a classical 2-DOF Quarter-Car suspension model. The model consists of two lumped masses:
\begin{itemize}
    \item \textbf{Sprung Mass ($m_s$)}: Represents one quarter of the vehicle chassis, cabin, and payload ($m_s = 300.0\,\text{kg}$).
    \item \textbf{Unsprung Mass ($m_u$)}: Represents the wheel rim, brake rotor, hub assembly, and tire mass ($m_u = 45.0\,\text{kg}$).
    \item \textbf{Suspension Elements}: Linear spring $k_s = 18,500\,\text{N/m}$ in parallel with viscous damper $c_s = 1,450\,\text{Ns/m}$ ($\zeta = 0.32$, $f_n = 2.50\,\text{Hz}$).
    \item \textbf{Tire Element}: Radial elastic spring $k_t = 190,000\,\text{N/m}$.
\end{itemize}

\begin{figure}[h!]
\centering
\begin{verbatim}
       +------------------------------------+
       |       Sprung Mass (ms = 300 kg)    | <--- z_s (Chassis Displacement)
       +------------------------------------+
                |                   |
            [Spring ks]        [Damper cs]
                |                   |
       +------------------------------------+
       |      Unsprung Mass (mu = 45 kg)    | <--- z_u (Axle Displacement)
       +------------------------------------+
                         |
                    [Tire kt]
                         |
       ====================================== <--- z_r(t) (Road Pothole Profile)
\end{verbatim}
\caption{Schematic of the 2-Degree-of-Freedom Quarter-Car Vehicle Dynamic Simulation Model.}
\label{fig:quarter_car_schematic}
\end{figure}

Applying Newton's second law to the sprung and unsprung masses yields the coupled second-order ordinary differential equations:
\begin{align}
    m_s \ddot{z}_s + c_s (\dot{z}_s - \dot{z}_u) + k_s (z_s - z_u) &= 0 \label{eq:ode_sprung} \\
    m_u \ddot{z}_u - c_s (\dot{z}_s - \dot{z}_u) - k_s (z_s - z_u) + k_t (z_u - z_r(t)) &= 0 \label{eq:ode_unsprung}
\end{align}

\subsubsection{State-Space Matrix Representation}
We formulate the system in continuous state-space form $\dot{\mathbf{x}}(t) = \mathbf{A}\mathbf{x}(t) + \mathbf{B} u(t)$, defining state vector $\mathbf{x} = [z_s, \, \dot{z}_s, \, z_u, \, \dot{z}_u]^T$ and scalar road excitation input $u(t) = z_r(t)$:

\begin{equation}
    \frac{d}{dt}
    \begin{bmatrix}
        z_s \\
        \dot{z}_s \\
        z_u \\
        \dot{z}_u
    \end{bmatrix}
    =
    \begin{bmatrix}
        0 & 1 & 0 & 0 \\
        -\frac{k_s}{m_s} & -\frac{c_s}{m_s} & \frac{k_s}{m_s} & \frac{c_s}{m_s} \\
        0 & 0 & 0 & 1 \\
        \frac{k_s}{m_u} & \frac{c_s}{m_u} & -\frac{k_s + k_t}{m_u} & -\frac{c_s}{m_u}
    \end{bmatrix}
    \begin{bmatrix}
        z_s \\
        \dot{z}_s \\
        z_u \\
        \dot{z}_u
    \end{bmatrix}
    +
    \begin{bmatrix}
        0 \\
        0 \\
        0 \\
        \frac{k_t}{m_u}
    \end{bmatrix}
    z_r(t)
    \label{eq:state_space_quarter_car}
\end{equation}

\subsubsection{Road Excitation Function $z_r(t)$}
When a vehicle navigates a pothole of longitudinal length $L_{\text{phys}}$ (meters) and maximum depth $d_{\max}$ (meters) at forward velocity $v$ (m/s), the tire traverses the depression over time interval $\Delta t_{\text{cavity}} = L_{\text{phys}} / v$. We model the road vertical profile as an inverted parabolic depression:
\begin{equation}
    z_r(t) = 
    \begin{cases}
        -4 \cdot d_{\max} \cdot \left(\frac{t - t_{\text{enter}}}{\Delta t_{\text{cavity}}}\right) \left(1 - \frac{t - t_{\text{enter}}}{\Delta t_{\text{cavity}}}\right), & t \in [t_{\text{enter}}, \, t_{\text{enter}} + \Delta t_{\text{cavity}}] \\
        0, & \text{otherwise}
    \end{cases}
    \label{eq:road_profile_ode}
\end{equation}

Equations (\ref{eq:state_space_quarter_car}) and (\ref{eq:road_profile_ode}) are solved numerically using explicit Runge-Kutta 4th/5th order integration (Dormand-Prince adaptive stepping with strict maximum step size $\Delta t_{\max} = 0.005\,\text{s}$). The resulting vertical chassis acceleration in units of gravity is:
\begin{equation}
    a_s(t) = \frac{\ddot{z}_s(t)}{9.81} = \frac{-c_s(\dot{z}_s - \dot{z}_u) - k_s(z_s - z_u)}{m_s \cdot 9.81} \quad [g]
\end{equation}

\begin{designreflection}{Suspension Isolation and Cabin Mount Transfer Functions}
During initial simulator testing, the quarter-car ODE solver predicted chassis shocks between $1.0g$ and $2.5g$ for moderate potholes. However, our empirical smartphone accelerometer dataset recorded median shocks of only $0.15g$ to $0.45g$ on identical road segments. 

I traced this discrepancy to vehicle cabin isolation. An accelerometer inside a vehicle cabin or mounted to a dashboard does not measure the bare suspension top-mount directly; it sits on a damped structural subframe and elastic dashboard plastic. Analyzing our IVP multi-mount dataset revealed that the cabin mount exhibits an empirical isolation attenuation factor $\alpha_{\text{mount}} = 0.3394 \approx 0.34$. 

By incorporating this empirical isolation factor ($a_{\text{phone}} = \alpha_{\text{mount}} \cdot a_s(t)$), the simulated peaks dropped to $0.159g$ for minor depressions and $0.417g$ for moderate potholes, aligning with real-world road shock distributions.
\end{designreflection}

\begin{designreflection}{Handling Adaptive Step Bounds in Dynamic Impact ODEs}
In numerical ordinary differential equation solving, high-stiffness spring-damper combinations such as automotive tires ($k_t = 190,000\,\text{N/m}$) create stiff numerical regimes. When I initially ran Scipy's \texttt{solve\_ivp} without step bounds, the adaptive Runge-Kutta solver took time steps larger than $0.05\,\text{s}$ during smooth road intervals, completely leaping across the brief $0.046\,\text{s}$ transient pothole traversal window without sampling any depression force! 

This resulted in trivial zero acceleration responses. To rectify this numerical artifact, I enforced a hard step-size constraint $\Delta t_{\max} = 0.005\,\text{s}$ (matching a $200\,\text{Hz}$ sampling frequency). This guaranteed that even at high velocities ($60\,\text{km/h}$), the numerical solver captures at least 8 to 12 quadrature points within the cavity impact zone, preserving sharp jerk discontinuities and wheel rebound transients.
\end{designreflection}

\subsection{Cross-Modal Consistency Validation and Severity Scoring}
To validate whether the predicted optical 3D geometry is physically coherent, the validation engine compares the simulated peak acceleration $\hat{a}_{\text{sim}}$ with the empirical expected shock envelope $a_{\text{exp}}$ calibrated from 96 ground-truth pothole waveforms:
\begin{equation}
    a_{\text{exp}}(d_{\max}, v) = 0.08 \cdot (d_{\max})^{0.82} \cdot \left(\frac{v}{35.0}\right)^{0.5} \quad [g]
\end{equation}

We evaluate the normalized physical consistency residual $\epsilon_{\text{res}}$:
\begin{equation}
    \epsilon_{\text{res}} = \frac{|\hat{a}_{\text{sim}} - a_{\text{exp}}|}{a_{\text{exp}} + 0.05}
\end{equation}

The continuous Dynamic Confidence Score $C_{\text{dyn}}$ is evaluated through a Gaussian radial basis function:
\begin{equation}
    C_{\text{dyn}} = \exp\left( - \frac{1}{2} \left(\frac{\epsilon_{\text{res}}}{0.5}\right)^2 \right) \in [0.0, \, 1.0]
    \label{eq:confidence_score_formula}
\end{equation}

A candidate detection is granted civil engineering validation if $C_{\text{dyn}} \ge 0.65$. If $\epsilon_{\text{res}}$ deviates significantly, the detection is flagged as either \texttt{SUSPICIOUS\_OVERESTIMATE} (suggesting optical glare or shadow artifacts exaggerating depth) or \texttt{SUSPICIOUS\_UNDERESTIMATE} (suggesting water-filled puddles masking true cavity depth).

\begin{table}[h!]
\centering
\small
\caption{End-to-End Pipeline Evaluation Census Across Sample Test Image Ingestion}
\label{tab:pipeline_evaluation_results}
\begin{tabular}{lcccccccl}
\hline
\textbf{Image Frame} & \textbf{Speed (km/h)} & \textbf{Width (cm)} & \textbf{Length (cm)} & \textbf{Depth (cm)} & \textbf{Sim Shock (g)} & \textbf{Confidence} & \textbf{Tier} & \textbf{Status} \\ \hline
\texttt{India\_000105} & 30.0 & 41.0 & 20.0 & 20.0 & 1.319 & 0.608 & Tier 3 & Suspicious Overestimate \\
\texttt{India\_000173} & 35.0 & 38.5 & 24.2 & 4.2 & 0.362 & 0.842 & Tier 2 & \textbf{PHYSICALLY VALIDATED} \\
\texttt{India\_000402} & 40.0 & 29.0 & 18.0 & 2.1 & 0.148 & 0.985 & Tier 1 & \textbf{PHYSICALLY VALIDATED} \\
\texttt{India\_000439} & 25.0 & 52.0 & 38.0 & 5.8 & 0.491 & 0.781 & Tier 2 & \textbf{PHYSICALLY VALIDATED} \\
\texttt{India\_000475} & 30.0 & 44.0 & 32.0 & 4.5 & 0.395 & 0.812 & Tier 2 & \textbf{PHYSICALLY VALIDATED} \\ \hline
\end{tabular}
\end{table}

The master inference pipeline integrates these modules into a single, cohesive executable ([`src/main_pipeline.py`](file:///d:/DSN3099-EPICS/Implementation/src/main_pipeline.py)), providing autonomous road infrastructure monitoring capable of processing raw dashcam video streams and producing verified 3D road condition audits in real time.
